% META: {"id": "TSPE-PROBA-EX-045", "chapitre": "TSPE-PROBABILITES", "type_objet": "exercice", "capacites": ["TSPE-PROBA-C9"], "capacites_codes": ["C9"], "methodes": ["M1"], "parcours": 1, "competences": ["calculer", "modeliser"], "duree_min": 8, "mode_creation": "ex_nihilo", "sources_inspiration": [], "fichier_tex": "chapitres/TSPE-PROBABILITES/exercices/TSPE-PROBA-EX-045.tex", "corrige_tex": "chapitres/TSPE-PROBABILITES/corriges/TSPE-PROBA-CO-045.tex", "status": "approved"}
% BEGIN-VERIFY
% from sympy import Rational
% n_min = 1 / (4 * Rational(5,100)**2 * Rational(1,10))
% assert n_min == 1000
% END-VERIFY
\begin{exercice}{TSPE-PROBA-EX-009}{2}{12}
On souhaite estimer une proportion $p$ inconnue dans une population, avec une precision $\delta=0{,}05$ et un risque $\alpha=0{,}1$ (c'est-a-dire $P(|M_n-p|\geq\delta)\leq\alpha$).
\begin{enumerate}
  \item Rappeler la majoration de $V(M_n)$ valable pour tout $p \in [0,1]$.
  \item En deduire une condition suffisante sur $n$ pour garantir $P(|M_n-p|\geq0{,}05)\leq0{,}1$.
  \item Calculer la plus petite valeur de $n$ verifiant cette condition.
\end{enumerate}
\end{exercice}
